\cotefile{134-2}
\scanpage{178}\mspage{33}
vérifier et d'expliciter comment les ``6 opérations'' habituelles sur les faisceaux (soit sur $X$, soit sur une réunion de strates de $X$) peuvent s'exprimer dans ce dictionnaire, dans le cas disons de strates non singulières (\unsure{autrement}, on aurait des difficultés avec les complexes dualisants, dont on aimerait qu'ils appartiennent à notre catégorie\unsure{.\,.\,.\,.\,.|\ldots}) \inserted{et de revérifier le formulaire connu de ces opérations}. Mais il \unsure{apparaitra|apparaîtra|apparaît} probablement que, pour bien pouvoir tout transcrire, il convient, plutôt \inserted{que} de regarder un diagramme du type
\diagram{\begin{tikzcd}[ampersand replacement=\&,column sep=1.4em,row sep=1.2em]
\cdot \arrow[r] \arrow[d] \& {} \\
{} \& {}
\end{tikzcd}}
dans la catégorie homotopique formée avec la cat. des ens. semi-simpliciaux, \inserted{de} regarder la catégorie des diagrammes d'ensembles $1/2$ simpliciaux, et à partir de là passer à une catégorie homotopique de fractions.

J'ai plus ou moins explicité dernièrement, en cogitant sur \deleted{\unsure{les|des|ces}}\unsure{des|les} fondements de ``topologie modérée'' (i.e. où on élimine d'emblée tous les phénomènes de sauvagerie), comment une stratification équisingulière\inserted{\deleted{(\unsure{d'un} \unsure{espace} \unsure{modéré} \unsure{compact})}}, à \unsure{strates} non singulières disons, d'un ``espace modéré'' \inserted{compact}, donne naissance canoniquement à un diagramme \deleted{\unsure{fon|f}} d'espaces,
% archive sheet number: 178
